\section*{Capítulo \ref{ch:g6:inside-the-earth} --- El interior de la Tierra}

\begin{solution}{exo:g6:inside-the-earth:1}
La \omterm{def:g6:inside-the-earth:layers}{corteza}, el \omterm{def:g6:inside-the-earth:layers}{manto}, el \omterm{def:g6:inside-the-earth:layers}{núcleo externo} y el \omterm{def:g6:inside-the-earth:layers}{núcleo interno}. El \omterm{def:g6:inside-the-earth:layers}{núcleo
externo} es líquido. Los \omterm{def:g6:inside-the-earth:layers}{núcleos externo} e interno están hechos de metal
(sobre todo hierro, con algo de níquel); la \omterm{def:g6:inside-the-earth:layers}{corteza} y el \omterm{def:g6:inside-the-earth:layers}{manto} están hechos
de \omterm{def:g3:stones-pebbles-sand:rock}{roca}.
\end{solution}

\begin{solution}{exo:g6:inside-the-earth:2}
(a) Falso: la \omterm{def:g6:inside-the-earth:layers}{corteza} es más gruesa bajo los \omterm{def:g3:continents-and-oceans:continent}{continentes} (unos 30~km) que
bajo los \omterm{def:g2:water-on-earth:ocean}{océanos} (unos 5~km). (b) Falso: el \omterm{def:g6:inside-the-earth:layers}{núcleo} está hecho de metal, una
aleación de hierro y níquel. (c) Verdadero. (d) Verdadero: 12\,262~m, unos
12.3~km.
\end{solution}

\begin{solution}{exo:g6:inside-the-earth:3}
(a) A 2891~km, unos 2900~km. (b) A 5150~km. (c) $5150 - 2891 = 2259$:
unos 2260~km.
\end{solution}

\begin{solution}{exo:g6:inside-the-earth:4}
``La cáscara es como la \underline{\omterm{def:g6:inside-the-earth:layers}{corteza}}, la clara es como el
\underline{\omterm{def:g6:inside-the-earth:layers}{manto}} y la yema es como el \underline{\omterm{def:g6:inside-the-earth:layers}{núcleo}}. Pero, a
diferencia de la yema, el \underline{\omterm{def:g6:inside-the-earth:layers}{núcleo}} tiene dos partes: un
\underline{\omterm{def:g6:inside-the-earth:layers}{núcleo externo}} líquido y un \underline{\omterm{def:g6:inside-the-earth:layers}{núcleo interno}}
sólido''.
\end{solution}

\begin{solution}{exo:g6:inside-the-earth:5}
$8.849 \div 6371 \approx 0.00139$, y $0.00139 \times 100 \approx 0.14$: el
Everest solo representa alrededor del \textbf{0.14~\%} del radio. Vista
desde lejos, la Tierra es una bola muy lisa: su \omterm{def:g1:landscapes-around-us:landform}{montaña} más alta es un
bultito diminuto.
\end{solution}

\begin{solution}{exo:g6:inside-the-earth:6}
(a) $30 \div 5 = 6$: seis veces más gruesa. (b) $100 \div 5 = 20$: veinte
veces más gruesa.
\end{solution}

\begin{solution}{exo:g6:inside-the-earth:7}
(a) $12 + 1 \times 25 = 37~^\circ$C. (b) $12 + 1.5 \times 25 = 12 + 37.5 =
49.5~^\circ$C. (c) $12 + 2 \times 25 = 62~^\circ$C.
\end{solution}

\begin{solution}{exo:g6:inside-the-earth:8}
La subida es de $74 - 14 = 60$ grados en 2~km, así que $60 \div 2 =
\textbf{30}$~grados por kilómetro. Sí: 30 está entre la media del este (25)
y la media del oeste (34).
\end{solution}

\begin{solution}{exo:g6:inside-the-earth:9}
(a) Justo encima del \omterm{def:g6:inside-the-earth:layers}{núcleo}, unos 13.7~km por segundo; justo por debajo de
su techo, unos 8.1~km por segundo: las ondas P se frenan de golpe.
(b) Entre 2891~km y 5150~km, en el \omterm{def:g6:inside-the-earth:layers}{núcleo externo}. (c) Las ondas S no
pueden atravesar un líquido: el \omterm{def:g6:inside-the-earth:layers}{núcleo externo} es líquido.
\end{solution}

\begin{solution}{exo:g6:inside-the-earth:10}
(a) $12.3 \div 30 = 0.41$: alrededor del \textbf{41~\%} de una \omterm{def:g6:inside-the-earth:layers}{corteza}
continental media. (b) No: bajo un \omterm{def:g3:continents-and-oceans:continent}{continente}, el \omterm{def:g6:inside-the-earth:layers}{manto} empieza por debajo
de la \omterm{def:g6:inside-the-earth:layers}{corteza}, a unos 30~km de profundidad de media; el sondeo se detuvo
mucho más arriba, dentro de la \omterm{def:g6:inside-the-earth:layers}{corteza}.
\end{solution}

\begin{solution}{exo:g6:inside-the-earth:11}
(a) $29 \times 6371 = 184\,759$, y $15 + 184\,759 = 184\,774$: unos
185\,000~grados. (b) El laboratorio da unos 6000~grados en el borde del
\omterm{def:g6:inside-the-earth:layers}{núcleo interno}: el número de Leo es unas 30~veces demasiado grande
($184\,774 \div 6000 \approx 31$). (c) El ritmo de 29~grados por kilómetro
solo se midió en los primeros 2~km, más o menos. Leo lo usa miles de
kilómetros más allá de las mediciones; en las profundidades, la
temperatura sube mucho más despacio. Un ritmo solo se puede usar allí
donde se midió.
\end{solution}

\begin{solution}{exo:g6:inside-the-earth:12}
(a) Si la \omterm{def:g6:inside-the-earth:layers}{corteza} es más ligera que la media, el material de más abajo
tiene que ser más denso que 5.5, para subir la media. (b) En el \omterm{def:g6:inside-the-earth:layers}{núcleo
externo} y el \omterm{def:g6:inside-the-earth:layers}{núcleo interno} (de unos 9.9 a unos 13~gramos por centímetro
cúbico); el \omterm{def:g6:inside-the-earth:layers}{manto} solo llega a 5.6 justo en su fondo. (c) Los materiales
más densos se hunden por debajo de los más ligeros, como una piedra se
hunde en el agua: el metal pesado acabó en el centro y las \omterm{def:g3:stones-pebbles-sand:rock}{rocas} más
ligeras, arriba.
\end{solution}

\begin{solution}{pb:g6:inside-the-earth:1}
\textbf{1.} $2 \times 6371 = 12\,742$~km. La maqueta mide 2~m de diámetro,
es decir, 2000~mm.

\textbf{2.} $12\,742 \div 2000 \approx 6.37$: 1~mm de la maqueta
representa unos 6.37~km de la Tierra real.

\textbf{3.} $8.849 \div 6.371 \approx 1.39$: el Everest mide unos
\textbf{1.4~mm} de altura en la maqueta, un grano de \omterm{def:g3:stones-pebbles-sand:sand}{arena}.

\textbf{4.} $12.3 \div 6.371 \approx 1.9$~mm. El centro de la maqueta está
1000~mm (1~m) por debajo de la superficie: el sondeo es unas 520~veces más
corto ($1000 \div 1.9 \approx 520$).

\textbf{5.} $2891 \div 6.371 \approx 453.8$: el techo del \omterm{def:g6:inside-the-earth:layers}{núcleo externo} se
pinta unos \textbf{454~mm} por debajo de la superficie.

\textbf{6.} $5150 \div 6.371 \approx 808.3$: unos 808~mm por debajo de la
superficie. El radio del \omterm{def:g6:inside-the-earth:layers}{núcleo interno} en la maqueta es de
$1000 - 808 = 192$~mm.

\textbf{7.} El \omterm{def:g6:inside-the-earth:layers}{manto}: de unos 5~mm a 454~mm, unos 449~mm. El \omterm{def:g6:inside-the-earth:layers}{núcleo externo}
ocupa $808 - 454 = 354$~mm, y el \omterm{def:g6:inside-the-earth:layers}{núcleo interno}, 192~mm.

\textbf{8.} $30 \div 6.371 \approx 4.7$~mm: una franja fina, que se puede
pintar con un pincel fino.

\textbf{9.} $15 + 2 \times 29 = 73~^\circ$C. En la maqueta, 2~km están
$2 \div 6.371 \approx 0.31$~mm por debajo de la superficie: apenas el
grosor de tres hojas de papel.

\textbf{10.} $4 \div 6.371 \approx 0.63$~mm, frente a unos 1.9~mm del
sondeo: la mina es menos profunda, por unos $1.9 - 0.6 = 1.3$~mm.

\textbf{11.} Que una sustancia se funda depende de la temperatura y de la
presión. El \omterm{def:g6:inside-the-earth:layers}{núcleo interno} es el más apretado de todos, por unos
3.3~millones de veces la presión del aire: el hierro apretado solo se funde
a una temperatura más alta, así que el hierro del \omterm{def:g6:inside-the-earth:layers}{núcleo interno} sigue
sólido aunque está más caliente que el \omterm{def:g6:inside-the-earth:layers}{núcleo externo}.

\textbf{12.} $5 \div 6.371 \approx 0.785$: la \omterm{def:g6:inside-the-earth:layers}{corteza} oceánica solo debe
pintarse con un espesor de \textbf{0.78~mm} en la maqueta de 2~metros, el
grosor de unas ocho hojas de papel ($0.78 \div 0.1 \approx 8$).
\end{solution}
